\originalchapter{顺磁体与磁热效应}\label{ch:magnetism}
\originalsection{磁矩的统计取向}
相互作用很弱的局域磁矩在外场中倾向降低能量, 热运动则倾向散布到所有方向或磁量子数. 设一个离子总角动量量子数为$J$, 磁矩在场方向的分量为$g\mu_Bm$, $m=-J,\ldots,J$, 含外场耦合的单离子能量为$\epsilon_m=-g\mu_BBm$. 用SI中的$B$及能量$-\vct\mu\cdot\vct B$约定, $\mu_B=e\hbar/(2m_e)$; 原课某些公式采用Gaussian单位, 不将其中$c$因子混入SI.

令$y=\beta g\mu_BB$, 几何级数给
\[
z_J=\sum_{m=-J}^J\e^{ym}
=\frac{\sinh[(J+1/2)y]}{\sinh(y/2)},\qquad
M=Ng\mu_B\partial_y\ln z_J.
\]
所以$M=Ng\mu_BJ B_J(y)$, 此处把Brillouin函数的自变量定义为$y$, 有
\[
B_J(y)=\frac1J\left[(J+1/2)\coth((J+1/2)y)-\frac12\coth(y/2)\right].
\]
部分教材把自变量换成$Jy$, 比较公式时需同步换元. 展开$\coth x=x^{-1}+x/3+O(x^3)$, 得弱场磁矩
\[
M\simeq\frac{N(g\mu_B)^2J(J+1)}{3\kb T}B,
\]
即Curie律. 强场低温时$M\to Ng\mu_BJ$, 因为最低能量态$m=J$占据; 饱和磁矩有限, 故不能把$1/T$弱场近似外推到$T=0$.

\originalsection{磁功的能量约定}
正则配分函数使用含Zeeman耦合的能量$E_Z=-BM$. 它的自由能$\mathcal F=-N\kb T\ln z_J$满足$\dd\mathcal F=-S\dd T-M\dd B$, 平均Zeeman能量为$E_Z=-BM$. 若用原课以总磁矩为外参的磁功$B\dd M$, 需定义场源Legendre转换后的内能$U=E_Z+BM$. 在理想无相互作用自旋模型中$U=0$, 并有$\dd U=T\dd S+B\dd M$. 这是两种能量记账之间的关系, 不意味着自旋能级没有场依赖.

熵$S=N\kb[\ln z_J-y\partial_y\ln z_J]$只依赖$B/T$. 对固定$B$加热, 交换热可写$C_B\dd T$, 其中$C_B=(\partial E_Z/\partial T)_B$; 对固定$M$的理想模型, $C_M=(\partial U/\partial T)_M=0$. 二者描述不同控制条件. 特别对$J=1/2$, 写单自旋磁矩大小$\mu=g\mu_B/2$, 则$z=2\cosh(\beta\mu B)$, $M=N\mu\tanh(\beta\mu B)$,
\[
E_Z=-N\mu B\tanh(\beta\mu B),\quad
C_B=N\kb(\beta\mu B)^2\operatorname{sech}^2(\beta\mu B).
\]
这个热容在两端都趋零, 中间有Schottky峰.

\originalsection{绝热去磁与负温}
磁矩系统可逆绝热时$S$固定, 非零熵的单调分支使$B/T$保持常数. 先在强场与热库接触使磁矩有序, 再隔热减小磁场, 会降低自旋温度. 实际晶格热容、残余磁相互作用和热漏阻止理想模型无限接近零温; 如果这些能量与外场能量可比, 仅用$B/T$不再足够.

在有上下能界的隔离自旋系统中, 更多磁矩占据高能级时可以有$\beta<0$. 此时高能态Boltzmann权重更大, 能量流向普通正温热库; 因而负温状态在热流意义上比所有正温状态都热. 需要分别说明自旋与晶格是否达到共同平衡. 把负温接到具有无上界平移动能的完整气体中通常无法维持一个共同负温正则系综.
